\begin{proposition}[The evaluated scalar and the two local leading terms]\hypertarget{tag-01RV}{}\label{inv:A2.2}
\StatementRecord{inv:A2.2}{proved}{01RV}{\SeriesRef{inv:eq:A-s-branch}, \SeriesRef{inv:eq:A-modular-K}, \SeriesRef{inv:A2.3}}

In the coordinate \eqref{inv:eq:A-s-branch},
\begin{equation}\hypertarget{tag-01T4}{}\label{inv:eq:A-mu-value}
\mu_s(0)=i\,C_4,
\qquad
C_4=\frac{\sqrt{2/3}}{\pi E_4(i)^{1/4}}
=\frac{4\sqrt\pi}{3^{3/4}\Gamma(1/4)^2}>0.
\end{equation}
Numerically,
\[
C_4=0.2366093140807852246877051134504496062943051340858243372192243072\ldots.
\]
For $r^3=t$, with $r>0$ on the arc $0<t<1$ and $\tau$ travelling
from $\rho$ to $i$ along the unit circle, set
\[
\lambda=\frac{3e^{5\pi i/6}}{2\pi E_6(\rho)^{1/3}}.
\]
Then $\tau=\rho+\lambda r+O(r^2)$ and
\[
K_r=-\frac{3i\lambda}{2\pi}r^2+O(r^3).
\]
The leading coefficient has real and imaginary parts
\[
\begin{aligned}
\operatorname{Re}K_2&=0.08010221108753827140763650309934451243\ldots,\\
\operatorname{Im}K_2&=0.13874109940222333968402480721280640441\ldots.
\end{aligned}
\]
\end{proposition}


\begin{proof}
\begin{enumerate}[label={$\langle1\rangle$\arabic*.},leftmargin=*]
\item Ramanujan's differential relations give
$E_6'(i)=-\pi i E_4(i)^2$ and
$E_4'(\rho)=-(2\pi i/3)E_6(\rho)$.
Using $j=1728E_4^3/(E_4^3-E_6^2)$, we obtain
\[
j(\tau)-1728=-1728\pi^2E_4(i)(\tau-i)^2+O((\tau-i)^3).
\]
Since $j=1728(1+s^4)$, the chosen branch is precisely
\eqref{inv:eq:A-s-branch}.  Formula \eqref{inv:eq:A-modular-K} now gives
\[
K_s(0)=-\frac4{\pi^2\sqrt{E_4(i)}}=6\mu_s(0)^2,
\]
because $\tau_s(0)=0$.  This proves nonvanishing and the first
formula in \eqref{inv:eq:A-mu-value}; the sign is fixed by the stated
coordinate convention.  Similarly
$j(\tau)=-512i\pi^3E_6(\rho)(\tau-\rho)^3+cdots$.
This gives the stated $\lambda$, and substitution into
\eqref{inv:eq:A-modular-K} gives the coefficient $-3i\lambda/(2\pi)$.

\item There is also a geometric nonvanishing proof independent of this
modular determinant.  Using the section identification proved in
\ref{inv:A2.3}, set $x=s^2X,y=s^3Y,t=1+s^4$.  The smooth family becomes
\[
Y^2=X^3-3(1+s^4)^3X+2(1+s^4)^4s^2,
\qquad P=(0,-\sqrt2(1+s^4)^2s)
\]
on the chosen branch.  At $s=0$ its equation is
$Y^2=X^3-3X$ and the point is $(0,0)$.  The first variation of the
curve itself is zero.  In the local coordinate $Y$ at that point,
$dX/(2Y)=dY/(3X^2-3)$, so the Abel integral of $P$ has derivative
$\sqrt2/3$.  The marked period is $-i\Omega_R$, rather than the
positive real period $\Omega_R$: the positive imaginary branch
$\tau=iy$, $y>1$, has elliptic scale $F/s=i/\Omega_R$ at zero.
Its period has zero first derivative.  Division gives
$\mu_s(0)=i\sqrt2/(3\Omega_R)\ne0$.
Using $g_2(\Z+i\Z)=(4\pi^4/3)E_4(i)$ and
$g_2(\Omega_R(\Z+i\Z))=12$ gives
$\Omega_R=\pi E_4(i)^{1/4}/\sqrt3$, which identifies this value
with \eqref{inv:eq:A-mu-value}.  This also explains why choosing the
positive real period with a positive real $s$ instead would select
the reciprocal branch $\operatorname{Im}\tau<1$.

\item The real period also has the integral expression
\[
\Omega_R=\frac{\sqrt2}{3^{1/4}}
\int_0^{\pi/2}(1-\tfrac12\sin^2\theta)^{-1/2}\,d\theta,
\qquad C_4=\frac{\sqrt2}{3\Omega_R}.
\]
It follows by the elliptic-integral substitution for the three roots
$\sqrt3,0,-\sqrt3$ of $Y^2=X^3-3X$ and the differential
$dX/(2Y)$. The beta-integral evaluation gives the gamma expression
in \eqref{inv:eq:A-mu-value}.

For the displayed numerical values,
\path{code/current_audits/audit_paper_i/code/a2_scalar_arb.py}
evaluates the divisor-sum series of $E_4(i)$ and $E_6(\rho)$
with 100 terms at 120 decimal working digits in Arb, and independently
evaluates the gamma expression for $C_4$.
The elementary inequalities $\sigma_3(n)\le n^4$,
$\sigma_5(n)\le n^6$, $e^{-2\pi}<1/500$ and
$e^{-\pi\sqrt3}<1/200$ give respective tail bounds
\[
250(N+1)^4 500^{-N-1},\qquad
1024(N+1)^6 200^{-N-1}\quad(N=100).
\]
The ratios of successive majorant terms are less than $1/480$
and $1/180$, respectively, which proves these tail estimates.
All arithmetic roundoff is included in the balls. The two $C_4$
enclosures overlap, and the resulting strict rational enclosures are
\[
\begin{aligned}
0.2366093140807852&<C_4<0.2366093140807853,\\
0.0801022110875382&<\operatorname{Re}K_2<0.0801022110875383,\\
0.1387410994022233&<\operatorname{Im}K_2<0.1387410994022234.
\end{aligned}
\]
All twelve checks, the narrower balls and their full inputs are
saved in \path{a2_scalar_arb.output.json}. The exact nonvanishing
and orders were proved in the preceding steps.
\end{enumerate}
\end{proof}